\subsection*{Performance by taxonomy}
Agent accuracy varied widely across functional systems, from olfaction and taste (11.2\%, mean over models), vision (11.3\%) and auditory and vestibular (13.5\%) to somatosensation (33.2\%). Within cross-system, evaluations on cortical and forebrain neurons (13.2\%, 10 evaluations) and on vasculature, barriers and CSF (14.3\%) were harder than those on glia and neuroimmune cells (16.7\%). This range should be read with caution, however. Most functional systems contain only three to nine evaluations, and pass rates vary as much within a system (pooled SD 17.9 percentage points across evaluations) as across the benchmark as a whole (17.8), so the differences between systems were not significant at the evaluation level (permutation test, $p = 0.50$) and are more likely to reflect which few evaluations happen to fall in each system than a property of the system itself.

By process class, agents were slightly more accurate on normal-function evaluations than on evaluations of disease, injury and aging (21.2\% vs. 19.0\%; Figure~\ref{fig:taxonomy}). Of the 21 models, 14 scored higher on normal function, as did 22 of the 34 configurations, but the gap was relatively small and not significant at the evaluation level (permutation test, $p = 0.53$).

By biological level, tissue- and organ-level evaluations were the hardest (13.0\%, 7 evaluations), molecular (19.7\%), cellular (19.7\%) and circuit (21.2\%) evaluations were intermediate, and behavioral evaluations were the easiest (25.1\%).
